\begin{lemma}
\label{lemma-finite-push-pseudo-coherent}
Let $A \to B$ be a ring map. Assume that $B$ is pseudo-coherent as an
$A$-module. Let $K^\bullet$ be a complex of $B$-modules.
The following are equivalent
\begin{enumerate}
\item $K^\bullet$ is $m$-pseudo-coherent
as a complex of $B$-modules, and
\item $K^\bullet$ is $m$-pseudo-coherent
as a complex of $A$-modules.
\end{enumerate}
The same equivalence holds for pseudo-coherence.
\end{lemma}

\begin{proof}
Assume (1). Choose a bounded complex of finite free $B$-modules
$E^\bullet$ and a map $\alpha : E^\bullet \to K^\bullet$ which is
an isomorphism on cohomology in degrees $> m$ and a surjection in degree $m$.
Consider the distinguished triangle
$(E^\bullet, K^\bullet, C(\alpha)^\bullet)$. By
Lemma \ref{lemma-recognize-pseudo-coherent}
$C(\alpha)^\bullet$ is $m$-pseudo-coherent as a complex of
$A$-modules. Hence it suffices to prove that $E^\bullet$ is
pseudo-coherent as a complex of $A$-modules, which follows from
Lemma \ref{lemma-complex-pseudo-coherent-modules}.
The pseudo-coherent case of (1) $\Rightarrow$ (2) follows from this and
Lemma \ref{lemma-pseudo-coherent}.

\medskip\noindent
Assume (2). Let $n$ be the largest integer such that $H^n(K^\bullet) \not = 0$.
We will prove that $K^\bullet$ is $m$-pseudo-coherent as a complex
of $B$-modules by induction on $n - m$. The case $n < m$ follows from
Lemma \ref{lemma-recognize-pseudo-coherent}.
Choose a bounded complex of finite free $A$-modules $E^\bullet$ and a
map $\alpha : E^\bullet \to K^\bullet$ which is an isomorphism on
cohomology in degrees $> m$ and a surjection in degree $m$.
Consider the induced map of complexes
$$
\alpha \otimes 1 : E^\bullet \otimes_A B \to K^\bullet.
$$
Note that $C(\alpha \otimes 1)^\bullet$ is acyclic in degrees
$\geq n$ as $H^n(E) \to H^n(E^\bullet \otimes_A B) \to H^n(K^\bullet)$
is surjective by construction and since $H^i(E^\bullet \otimes_A B) = 0$
for $i > n$ by the spectral sequence of
Example \ref{example-tor}.
On the other hand, $C(\alpha \otimes 1)^\bullet$
is $m$-pseudo-coherent as a complex of $A$-modules because
both $K^\bullet$ and $E^\bullet \otimes_A B$ (see
Lemma \ref{lemma-complex-pseudo-coherent-modules})
are so, see
Lemma \ref{lemma-cone-pseudo-coherent}.
Hence by induction we see that $C(\alpha \otimes 1)^\bullet$
is $m$-pseudo-coherent as a complex of $B$-modules. Finally
another application of
Lemma \ref{lemma-cone-pseudo-coherent}
shows that $K^\bullet$ is $m$-pseudo-coherent as a complex of $B$-modules
(as clearly $E^\bullet \otimes_A B$ is pseudo-coherent as a complex
of $B$-modules). The pseudo-coherent case
of (2) $\Rightarrow$ (1) follows from this and
Lemma \ref{lemma-pseudo-coherent}.
\end{proof}